\documentclass[10pt,a4paper]{amsart}
\usepackage[margin=19mm]{geometry}
\usepackage{amsmath,amssymb,mathtools}
\usepackage[scr=rsfs,cal=euler]{mathalfa}
\usepackage{xfrac}
\usepackage[unicode,hidelinks,bookmarks=false]{hyperref}
\newcommand{\ie}{i.e.\ }
\newcommand{\cf}{cf.\ }
\newcommand{\resp}{respectively\ }
\newcommand{\Subsection}{\subsection}
\providecommand{\nobreakdash}{\nobreak\hbox{-}}
\setlength{\emergencystretch}{2em}
\multlinegap0pt
\usepackage{bm}
\usepackage{bbm}
\usepackage{dsfont}
\usepackage{xspace}
\usepackage{cancel}
\usepackage{mathtools}
\usepackage{amsthm}
\makeatletter
\@ifundefined{theorem}{\newtheorem{theorem}{Theorem}}{}
\@ifundefined{lemma}{\newtheorem{lemma}[theorem]{Lemma}}{}
\makeatother
\def\cC{\mathcal{C}}
\def\cN{\mathcal{N}}
\def\dv{\,\mathrm{dv}}
\def\d{\,\mathrm{d}}
\def\idm{\int_{\partial M}}
\def\im{\int_{M}}
\def\r{\mathbb{R}}
\title[Sign-changing solutions to the Yamabe problem on manifolds w]{Sign-changing solutions to the Yamabe problem on manifolds with boundary}
\author{M\'onica Clapp, Benedetta Pellacci, Angela Pistoia}
\date{Source: arXiv:2511.10553v2 (CC BY 4.0)}
\begin{document}
\maketitle

\noindent\emph{Source and licence.} Sign-changing solutions to the Yamabe problem on manifolds with boundary. Version arXiv:2511.10553v2 (CC BY 4.0), distributed under the Creative Commons Attribution 4.0 International licence, as recorded in the arXiv licence metadata for this exact version and shown on the abstract page https://arxiv.org/abs/2511.10553v2. Full rights notice: licenses/arxiv-2511.10553-CC-BY-4.0.txt.\par\smallskip

\noindent\textit{Editorial scope.} Counted unit: the Lemma whose source label is \texttt{lem:nehari} in the article (source0.tex lines 256--262, source label \texttt{lem:nehari}), delivered together with the article's own complete original proof of it (source0.tex lines 264--279, one \texttt{proof} environment of 16 lines). No mathematics is written, repaired or re-derived: statement and proof are the article's own text line for line. Required context retained uncounted: eq:Q\_norm (lines 233--235). Every definition, display and label of those lines is retained unchanged. Omitted from this package and not redistributed: 1--232, 236--255, 263--263, 280--835. Nothing inside a retained excerpt is omitted or altered: every retained body is delivered exactly as the article's lines are written, so no omission receipt of any kind accompanies this package. Citations: the retained excerpts carry no citation command at all, so no bibliography entry of the article is transcribed and no reference is redistributed. Reference adaptation: none was needed and none was made; every reference inside the delivered unit resolves inside it, so no unresolved reference or citation is produced. Printed slips kept literal: no printed word, symbol, hypothesis or number of the counted statement or of its proof was altered. Layout and class: the article's own class is \texttt{amsart} with packages including inputenc, amsmath, amsfonts, amssymb, graphicx, verbatim, mathrsfs, upref, amsthm, amsxtra, exscale, cite, dsfont, geometry; the wrapper is the lane's pinned \texttt{amsart} editorial wrapper at 10pt on A4, which restates the theorem-like environments the retained text needs and adds the article's own macro definitions that the retained text needs (\texttt{\textbackslash cC}, \texttt{\textbackslash cN}, \texttt{\textbackslash d}, \texttt{\textbackslash dv}, \texttt{\textbackslash idm}, \texttt{\textbackslash im}, \texttt{\textbackslash r}), with extra wrapper packages (bm, bbm, dsfont, xspace, cancel, mathtools). The delivered numbering is the unit's own. Assets: no retained span contains a figure environment, a graphics-inclusion command, a PDF-page inclusion, a table or a TikZ picture; no image file is copied into this package, the package declares no asset dependency and no image file is redistributed with it. Licence: version arXiv:2511.10553v2 of this article is distributed under the Creative Commons Attribution 4.0 International licence, as recorded in the arXiv licence metadata for this exact version and shown on the abstract page https://arxiv.org/abs/2511.10553v2. A full rights notice is delivered at licenses/arxiv-2511.10553-CC-BY-4.0.txt. Characters: every retained excerpt is pure ASCII, so the delivered engine, XeLaTeX with its default Latin Modern text font, reports no missing character.
\begin{equation} \label{eq:Q_norm}
C_1\|u\|^2 \leq \mathscr{Q}u \leq C_2\|u\|^2\qquad\text{for all \ }u\in H^1(M).
\end{equation}

\begin{lemma}\label{lem:nehari}
\begin{itemize}
\item[$(a)$]There exists $c_0>0$ such that $\|u\|\geq c_0$ for all $u\in\cN$.
\item[$(b)$]$\cN$ is a Hilbert submanifold of class $\cC^2$ of $H^1(M)$ and a natural constraint for $J$.
\item[$(c)$]If $u\in H^1(M)\smallsetminus\{0\}$, then there exists a unique $t_u\in(0,\infty)$ such that $t_uu\in\cN$. The function $t\mapsto J(tu)$ is strictly increasing in $[0,t_u]$ and strictly decreasing in $[t_u,\infty)$.
\end{itemize}
\end{lemma}

\begin{proof}
$(a):$ \ If $u\in\cN$, then, by \eqref{eq:Q_norm} and the Sobolev embeddings, there is $c_1>0$ such that
$$\mathscr{Q}u=a_n\im |u|^p \dv + b\idm|u|^q\d\sigma\leq c_1\Big((\mathscr{Q}u)^{p/2}+(\mathscr{Q}u)^{q/2}\Big).$$
Therefore, $\mathscr{Q}u\geq c_0>0$ for all $u\in\cN$, and the statement $(a)$ follows from \eqref{eq:Q_norm}.

$(b):$ \ It follows from $(a)$ that $\cN$ is a closed subset of $H^1(M)$. The function $\Psi(u):=\mathscr{Q} u-F(u)$ is of class $\cC^2$ on $H^1(M)$ and $\cN=\Psi^{-1}(0)\smallsetminus\{0\}$. Since
\begin{align}\label{eq:psi}
\Psi'(u)u&=2\mathscr{Q}u-F'(u)u=2\mathscr{Q}u-pa_n\im |u|^p\dv-qb\idm |u|^q\d \sigma \nonumber\\
&\leq 2\mathscr{Q}u-qF(u) =(2-q)\mathscr{Q}u<0\qquad\text{for every \ }u\in\cN,
\end{align}
$0$ is a regular value of $\Psi$. This shows that $\cN$ is a Hilbert submanifold of class $\cC^2$ of $H^1(M)$. Furthermore, \eqref{eq:psi} shows that $u$ does not belong to the tangent space $T_u\cN$ of $\cN$ at $u$. Hence, $H^1(M)=T_u\cN\oplus\r u$. So, if $u\in\cN$ and $J'(u)v=0$ for every $v\in T_u\cN$ then, as $J'(u)u=0$, we have that $J'(u)v=0$ for every $v\in H^1(M)$. This shows that every critical point of the restriction of $J$ to $\cN$ is a critical point of $J$, as claimed.

$(c):$ \ Let $u\in H^1(M)\smallsetminus\{0\}$ and $J_u:[0,\infty)\to\r$ be given by
$$J_u(t):=J(tu)=\frac{\mathscr{Q}u}{2} t^2- a_ut^p-b_ut^q,\quad\text{with \ }a_u:=\frac{a_n}{p}\im|u|^p\dv, \ b_u:=\frac{b}{q}\idm|u|^q\d\sigma.$$
As $\mathscr{Q}u>0$ and $a_u>0$, $J_u$ has a unique critical point $t_u$ in $(0,\infty)$ which is a global maximum. Note that, if $t\in(0,\infty)$, then $tu\in\cN$ iff $t_u$ is a critical point of $J_u$. This completes the proof.
\end{proof}

\end{document}
